  
\addcontentsline{toc}{chapter}{\twkkai{对话录·其一}}  
\markboth{\twkkai{对话录·其一}}{}  
\pagecolor{black}
\color{white}
\quad\quad \twkkai{\lettrine[,lraise=0.2,lhang=0.5]{\textbf{猫}}{：}开开门呗。

守门人\index{shou@守门人}：……需要通行奖章。

猫：别这样嘛——我的奖章被我弄丢了。

守门人：……这不在我的职责范围内。

猫：我知道暗语。“狄利克雷\index{dilikelei@狄利克雷（P.\,G.\,L.\ Dirichlet）}”。

守门人：……进来吧。

……

猫：哎哟，我的朋友，可折腾死我了！你猜怎么着——使者没找着，反倒落下了我的奖章。

狄利克雷：那你怎么回来的？奖章可是你的导航仪啊。

猫：幸好我没走太远——仅此而已。不过你想听听我的见闻吗？

狄利克雷：噢？请讲。

猫：我第一次听见人间的乐声——虽然和我们平时听的相比，水平还是差了点；不过倒是挺让我惊讶的。它们和以你的名字命名的特征\(\chi(n)\)可有点关系呢。

狄利克雷：我猜你并不是想说我发现的特征，而是傅里叶变换吧。

猫：啊——是的。你不是写出过\(L\)-函数的定义式吗？

狄利克雷：是啊。它的定义式是……

守门人：……\(L(s,\chi)=\sum_{n=1}^{\infty}\frac{\chi(n)}{n^s}=\prod_{p}\left(1-\frac{\chi(p)}{p^s}\right)^{-1},\  \Re(s)>1.\)

狄利克雷：守门人在呢喃什么呢？怎么抢了我的话……

猫：不用理他；准是又在打瞌睡说梦话了。不过他倒是梦呓到点上了。

狄利克雷：嗯。对。那这条式子怎么了呢？

猫：它不就是你所说的傅里叶变换\index{fu@傅里叶变换（与 $L$-函数的类比）}里的时域性\index{fu@傅里叶变换（与 $L$-函数的类比）!时域性}的完美呈现吗？具体来说，每个正整数\(n\)对应有一个振幅\index{fu@傅里叶变换（与 $L$-函数的类比）!振幅}——\(\chi(n)/n^s\)。我们可爱的\(\sum\)便把这其中的基频\index{fu@傅里叶变换（与 $L$-函数的类比）!基频}（暂且把\(n=1\)对应的振幅称为基频吧）和泛音\index{fu@傅里叶变换（与 $L$-函数的类比）!泛音}全部相加，得到一个\(L\)-函数。

狄利克雷：看你平时的样子——你这次不会又是胡说的吧；虽然你这次讲得头头是道。

猫：我保证绝对不是你所想的那样的！仔细想想，一个\(L\)-函数由其特征\(\chi\)所决定，成为一种独特的乐器；但再想——世上的\(\chi\)有无数种，从而就衍生出了无数种乐器！它们就构成了我听到的协奏曲\index{xiezouqu@协奏曲（Concerto）}！

狄利克雷：好啦……说了那么多，你不打算给我了解了解你的见闻吗？

猫：正有此意。不过奖章不见了，现在没有办法。我倒有一计——给你跳段舞如何\marginnote{
\centering
{\Large ♪}

\vspace{0.4em}

\qraudio{Allegro}{\AllegroSound}{https://raw.githubusercontent.com/kelvinlamresearchmath-web/finding-dirichlet-site/master/Allegro.mp3}

\vspace{0.3em}
}[1.6cm]
{\small ♪}？

狄利克雷：来吧。

……

\noindent\adjustbox{center}{\includegraphics[scale=0.7]{Allegro.pdf}}

狄利克雷：看你扭得真辛苦啊——猫爪都快磨没了吧；发出的声音倒是美妙。

猫：哎哟——此曲只应天上有，说得可没错了。你倒是要好好珍惜这宝贵的音乐啊——爪子可以再长回来，但曲子可不是每时每刻都有的。

狄利克雷：嗯。那这是谁作的曲呢？

猫：你一定认识他。

狄利克雷：噢？

猫：他是时代的天才——莫扎特\index{mozhate@莫扎特（W.\,A.\ Mozart）}呢！刚才这段，是莫扎特K.622第一乐章的第一主题！而且莫扎特和你或许也有点联系。

狄利克雷：愿闻其详。

猫：莫扎特的妻子是作曲家韦伯\index{wei@韦伯（C.\,M.\ von Weber）}的堂妹；韦伯又认识门德尔松\index{men@门德尔松（F.\ Mendelssohn）}。门德尔松的妹妹不就是你的妻子\footnote{莫扎特与康斯坦茨·韦伯结婚；康斯坦茨与作曲家卡尔·玛利亚·冯·韦伯为堂兄妹；韦伯与费利克斯·门德尔松同属早期浪漫派，交游密切；费利克斯之妹雷贝卡·门德尔松嫁给数学家狄利克雷。}吗？

狄利克雷：这么巧！不过我怀疑你是杜撰的。

猫：任你怎么怀疑好了。反正，这首曲子是莫扎特为他的挚友——一位单簧管\index{danhuangguan@单簧管（clarinet）}演奏家——写的。当然，刚才的乐句在小提琴部分就已经以某种类似的方式先演绎过了——就是说，开头的几个音符构成的动机（\textit{motif}）\index{dong@动机（motif）}是相同的；但后面略有变动。这就是我所说的多种\(L\)-函数组合而成协奏曲的生动体现。

狄利克雷：你还是跳舞吧。

猫：可别嫌弃我所说的话啊。不过我倒乐意为你效劳\marginnote{
\centering
{\Large ♪}

\vspace{0.4em}

\qraudio{Adagio}{\AdagioSound}{https://raw.githubusercontent.com/kelvinlamresearchmath-web/finding-dirichlet-site/master/Adagio.mp3}

\vspace{0.3em}
}[0.9cm]
{\small ♪}。

……

\noindent\adjustbox{center}{\includegraphics[scale=0.7]{Adagio.pdf}}

狄利克雷：怎么变风格了？听得我犯困了。

猫：这段来自莫扎特K.622第二乐章。\textit{Adagio}——柔板，自然是这样的。你知道，作曲家们写古典协奏曲时都喜欢快-慢-快形式的写法\index{xiezouqu@协奏曲（Concerto）!快-慢-快乐章结构}。

狄利克雷：噢。了解了。现在我倒宁愿你继续讲回你的理论了。

猫：喵，好的。我的朋友，你知道\(L\)-函数可以被分解成乘积形式的吧？正如守门人的梦呓那般。

狄利克雷：嗯。接着呢？

猫：正如乐曲演奏需要乐谱那般，如若需要发出各种音色，就需要不同的生成器。于是你便可以把\((1-\chi(p)p^{-s})^{-1}\)这个组件想象成一个生成器。把它拆开成\(1+\chi(p)p^{-s}+\chi(p^2)p^{-2s}+\cdots\)，你就得到第\(p\)号素泛音及其所有倍频（\(p^2,p^3,\dots\)）的振幅了。这体现的是傅里叶变换中的频域性\index{fu@傅里叶变换（与 $L$-函数的类比）!频域性}。

狄利克雷：我的内心很兴奋，以至于我必须和你说出来——我似乎理解你话里的那奥妙了！

猫：嘿，是吧！那允许我再舞一曲吧\marginnote{
\centering
{\Large ♪}

\vspace{0.4em}

\qraudio{Rondo}{\RondoSound}{https://raw.githubusercontent.com/kelvinlamresearchmath-web/finding-dirichlet-site/master/Rondo.mp3}

\vspace{0.3em}
}[1.6cm]
{\small ♪}。

狄利克雷：请便。

……

\noindent\adjustbox{center}{\includegraphics[scale=0.7]{Rondo.pdf}}

狄利克雷：这段也太快了。你不累吗？

猫：是挺累的，这段是莫扎特K.622第三乐章节选，是\textit{Rondo: Allegro}呢。哎，这里似乎越来越热了。噢！忘了我的奖章了——它可以让所有这一切变得更简单。

狄利克雷：既然话都说到这里了，我觉得你还是尽快出去再找找好一点。那奖章可背负着重大的使命呢。我也累了；你快去快回吧——希望你有新的见闻。

猫：喵。那我先去找奖章了；你先睡吧。

\rule{0pt}{10pt}

\noindent\adjustbox{center}{\includegraphics[scale=0.7]{trebleclef.preview.pdf}}

}


